\section{Big jumps and the boundary layer}\label{sec:bigjump}

Outside the bulk each bin is $o(1/B_N)$ by Theorem~\ref{thm:bulk}. How much
smaller it is depends on its distance from $pN$. On the ballistic range
$k-pN\asymp N$ the bins are of order $T_\plus(N)$, smaller than $1/B_N$ by a
power of $N$. Near $k-pN=KB_N$ with $K$ fixed they are still of order $1/B_N$,
with a factor of order $K^{-\alpha_\plus-1}$ when the plus scale is the larger one
(Remark~\ref{rem:no-single-formula}). This section describes these bins and the
boundary layer with one kernel.

Throughout we write $j=N-k$ for the number of minus slots. For $k>pN$ put
\begin{equation}\label{eq:Lk}
 L_k=\frac{k-pN}{1-p},\qquad\text{so that}\qquad N-L_k=\frac{j}{1-p}=\frac{js}{\mu_\minus}.
\end{equation}
This is the length of the plus run that moves a typical configuration to the
plus count $k$. For a start of sign $\sgn_0$ and $j\ge1$ put
\begin{equation}\label{eq:gamma}
 \gamma^{(\sgn_0)}_\plus(j)=\sum_{i=0}^{j-1}\kappa^{(\sgn_0)}_\plus(i)\bigl(1+\mu_\plus u_\minus(j-i)\bigr),
 \qquad\gamma^{(\sgn_0)}_\plus(0)=0,
\end{equation}
and define the big-jump kernel by
\begin{equation}\label{eq:G}
 G^{(\sgn_0)}_N(k)=\gamma^{(\sgn_0)}_\plus(N-k)\,f_\plus(\lfloor L_k\rfloor)
 +T_\plus(\lfloor L_k\rfloor)\,\kappa^{(\sgn_0)}_\plus(N-k),\qquad pN<k\le N .
\end{equation}
The first term is a complete plus run of length about $L_k$. The coefficient
$\gamma^{(\sgn_0)}_\plus(j)$ counts its admissible positions. The run starts right
after the $i$-th minus slot, which has probability $\kappa^{(\sgn_0)}_\plus(i)$, and
the fresh minus walk that follows it must hold the other $j-i$ minus slots. By the
Green identity, Lemma~\ref{lem:green}, the expected number of lengths of that walk
that do so is $1+\mu_\plus u_\minus(j-i)$. The second term is a censored final
plus run of length at least $L_k$. By~\eqref{eq:renewal-thm} and~\eqref{eq:kappa},
$\kappa^{(\sgn_0)}_\plus(i)\to1/\mu_\minus$ and $u_\minus(i)\to1/\mu_\minus$, so a
Ces\`aro average gives
\begin{equation}\label{eq:gamma-asymp}
 \gamma^{(\sgn_0)}_\plus(j)\sim\frac{js}{\mu_\minus^2}=\frac{N-L_k}{\mu_\minus}\qquad(j\to\infty).
\end{equation}
At $k=N-j$ with $j$ fixed the first term is of order $T_\plus(N)/N$ and the
second is $T_\plus(N)\kappa^{(\sgn_0)}_\plus(j)$ up to a factor $1+o(1)$, so the
kernel reduces to the boundary layer. As $j\to\infty$ the coefficient of the
second term tends to $1/\mu_\minus$, its value in the interior. So one kernel
covers the big-jump zone and the boundary layer.

The zone starts at a multiple of $B_N$ when the plus scale is the larger one.
Otherwise it starts a little later, and we put
\begin{equation}\label{eq:lambdaN}
 \lambda_N=1+\log\frac{B_N}{b_{\plus,N}}=1+\log\frac1{\theta_{\plus,N}}\ \ge1 .
\end{equation}
So $\lambda_N=1$ when $b_{\plus,N}=B_N$, for instance for all large $N$ when
$\alpha_\plus<\alpha_\minus$, and $\lambda_N=O(\log N)$ in all cases.

\begin{introtheorem}[big jumps and the boundary layer]\label{thm:bigjump}
Let the start be fresh with sign $\sgn_0$, and write $\kappa=\kappa^{(\sgn_0)}_\plus$
and $G_N=G^{(\sgn_0)}_N$.
\begin{enumerate}[(a)]
\item (The boundary layer.) Assume (RV) for the plus sign. For every fixed $j\ge0$,
\[
 \frac{P_N(N-j)}{T_\plus(N)}\to\kappa(j).
\]
\item (The big-jump zone.) Assume (RV) and (SPAN) for both signs and (LR) for the
plus sign. Then
\[
 \lim_{K\to\infty}\ \limsup_{N\to\infty}\ \sup_{pN+K\lambda_NB_N\le k\le N}
 \Bigl|\frac{P_N(k)}{G_N(k)}-1\Bigr|=0,
\]
where $G_N(k)=0$ only if $P_N(k)=0$, and then the ratio is read as $1$. In
particular, with the size function
\begin{equation}\label{eq:Ghat}
 \widehat G_N(k)=\frac N{L_k}\,T_\plus(\lfloor L_k\rfloor),
\end{equation}
we have
\[
 \lim_{K\to\infty}\ \limsup_{N\to\infty}\ \sup_{pN+K\lambda_NB_N\le k\le N}
 \frac{|P_N(k)-G_N(k)|}{\widehat G_N(k)}=0,
\]
and this additive form also holds for the kernel $\overline G_N=\overline G^{(\sgn_0)}_N$
obtained from~\eqref{eq:G} by replacing $\gamma^{(\sgn_0)}_\plus(N-k)$ with
$(N-L_k)/\mu_\minus$.
\item (Uniform boundary layer.) Under the hypotheses of (b), for every
$\delta_N\to0$,
\[
 \sup_{0\le j\le\delta_NN}\Bigl|\frac{P_N(N-j)}{T_\plus(N)}-\kappa(j)\Bigr|\to0
\]
and
\[
 \sup\Bigl\{\Bigl|\frac{P_N(N-j)}{T_\plus(N)\kappa(j)}-1\Bigr|\ \Big|\ 0\le j\le\delta_NN,\ \kappa(j)>0\Bigr\}\to0 .
\]
\end{enumerate}
The same holds at the minus end with the signs exchanged.
\end{introtheorem}

The exact coefficient $\gamma^{(\sgn_0)}_\plus(j)$ matters for the ratio form at
bounded $j$ with $\kappa(j)=0$. If $f_\minus(1)=0$, then for a plus start
$P_N(N-1)=f_\plus(N-1)$ and $G_N(N-1)=f_\plus(\lfloor N-s/\mu_\minus\rfloor)$, whose ratio
tends to one by (LR), while the limit coefficient $s/\mu_\minus^2$ would not. In the
additive form the two coefficients agree. In (c) the condition $\kappa(j)>0$ excludes
nothing but $j=0$ for a minus start when $f_\minus(1)>0$.

The double limit in (b) is forced. Fix $K$ and let $k=pN+KB_N$.
Theorem~\ref{thm:bulk}(c) gives $B_NP_N(k)=h_N(x_{N,k})+o(1)$. By
Remark~\ref{rem:no-single-formula} the ratio of $h_N(K)$ to $B_NG_N(k)$ tends to one
only as $K\to\infty$. When $\theta_{\plus,N}\to0$ a factor in front of $B_N$ that
tends to infinity is also needed (Remark~\ref{rem:theta}). Then on the window
$KB_N\le k-pN\le K\lambda_NB_N$ we have only the additive bound of
Theorem~\ref{thm:bulk}(a) and the lower bound $P_N(k)\ge(1-o(1))G_N(k)$ of
Corollary~\ref{cor:bigjump-lower}, and we do not describe $P_N$ there. This window
lies on the side of the sign with the smaller scale. When the minus scale is the
smaller one it is the window at the minus end given by the same statements with
the signs exchanged.

\subsection{Giant runs}\label{sec:giant}

Every bin in this section is decomposed according to the plus runs that occupy
many slots. We use the weights $\rho^{(\sgn_0)}_i(x)$ of
Section~\ref{sec:kernels}, with the representation~\eqref{eq:rho-rep}, and for
$i\ge1$ and $x\ge0$ we put
\[
 \beta_i(x)=\Pi^\minus_{i+x}(i),
\]
the probability that a fresh minus walk on $i+x$ slots has exactly $i$ minus
slots. Then
\begin{equation}\label{eq:weight-masses}
 \sum_{x\ge0}\rho^{(\sgn_0)}_i(x)=\kappa^{(\sgn_0)}_\plus(i),\qquad
 \sum_{x\ge0}\beta_i(x)=1+\mu_\plus u_\minus(i),
\end{equation}
by~\eqref{eq:rho-mass} and by the Green identity, Lemma~\ref{lem:green}.

We need $\beta_i$ in terms of sums of runs. Take a fresh minus walk with minus
runs $X^\minus_1,X^\minus_2,\dots$ and plus runs $X^\plus_1,X^\plus_2,\dots$ in the
order $X^\minus_1X^\plus_1X^\minus_2X^\plus_2\cdots$. Let $J=J_\minus(i)$ be the
index of the minus run that contains the $i$-th minus slot, in the notation of
(K3) in Appendix~\ref{sec:bulk-kernels}, and put
$\xi_i=X^\plus_J\ones_{\{S^\minus_J=i\}}$. The $i$-th minus slot is slot
$i+S^\plus_{J-1}$. The next minus slot follows it at once unless the run
$X^\minus_J$ ends there, and then it follows after $X^\plus_J$ plus slots. Hence
\begin{equation}\label{eq:beta-rep}
 \beta_i(x)=\Prob\bigl(S^\plus_{J-1}\le x\le S^\plus_{J-1}+\xi_i\bigr).
\end{equation}
Here $J$ depends only on the minus runs, and $X^\plus_J$ is independent of
$(J,S^\plus_{J-1},X^\minus_1,\dots,X^\minus_J)$.

For an integer $g\ge1$ let $\mathcal G_g$ be the number of plus runs that have at
least $g$ slots among the first $N$ slots. We call these runs giant. A complete run
counts with its length, and the final run counts with the number of its slots up to
slot $N$. Let $M_N$ be the largest number of slots that a plus run has among the
first $N$ slots, so $\{M_N<g\}=\{\mathcal G_g=0\}$.

\begin{lemma}[the giant identity]\label{lem:giant}
For every start of sign $\sgn_0$, every $N$, every $0\le j<N$ with $k=N-j$, and
every integer $g\ge1$,
\begin{equation}\label{eq:giant}
 \E\bigl[\mathcal G_g\ones_{\{A_N=k\}}\bigr]=\mathrm{Cen}_g(k)+\mathrm{Int}_g(k),
\end{equation}
where, with $\rho_i=\rho^{(\sgn_0)}_i$,
\begin{align*}
 \mathrm{Cen}_g(k)&=\sum_{x\ge0,\ N-j-x\ge g}\rho_j(x)\,T_\plus(N-j-x),\\
 \mathrm{Int}_g(k)&=\sum_{i=0}^{j-1}\ \sum_{x,x'\ge0,\ N-j-x-x'\ge g}\rho_i(x)\,\beta_{j-i}(x')\,f_\plus(N-j-x-x').
\end{align*}
Consequently
\begin{equation}\label{eq:giant-split}
 P_N(k)=\Prob(A_N=k,\ M_N<g)+\mathrm{Cen}_g(k)+\mathrm{Int}_g(k)
 -\E\bigl[(\mathcal G_g-1)^+\ones_{\{A_N=k\}}\bigr],
\end{equation}
and, with $f^*_\plus(n)=\sup_{m\ge n}f_\plus(m)$,
\begin{equation}\label{eq:giant-crude}
 \mathrm{Cen}_g(k)\le\kappa^{(\sgn_0)}_\plus(j)\,T_\plus(g),\qquad
 \mathrm{Int}_g(k)\le\gamma^{(\sgn_0)}_\plus(j)\,f^*_\plus(g).
\end{equation}
\end{lemma}

\begin{proof}
By linearity, the left side of~\eqref{eq:giant} is the sum over all plus runs of
the probability that the run is giant and $A_N=k$. We split the runs into the
final run and the complete runs.

Let the final run be a plus run that starts at slot $t+1$. All $j$ minus slots
precede it, so $t=j+x$, where $x$ is the number of plus slots before it. The run
has $N-j-x$ slots up to slot $N$, and its length is at least $N-j-x$. The length is
independent of the runs before it. So this case contributes
$\rho_j(x)T_\plus(N-j-x)$ for each $x$ with $N-j-x\ge g$, which is
$\mathrm{Cen}_g(k)$.

Let a complete plus run start right after the $i$-th minus slot, with $x$ plus
slots before it, and let its length be $y$. It occupies slots $i+x+1$ to $i+x+y$,
and $i+x+y\le N-1$. The runs after it form a fresh minus walk, independent of
everything before. On $\{A_N=k\}$ this walk holds the remaining $j-i$ minus slots
among its $N-i-x-y\ge1$ slots. Since its first slot is a minus slot, $j-i\ge1$.
Write $N-i-x-y=(j-i)+x'$. Then $y=N-j-x-x'$, and this case contributes
$\rho_i(x)f_\plus(y)\beta_{j-i}(x')$. A fresh plus start has $\rho_0=\delta_0$, and
then the first run has law $f_\plus$, as every other plus run. Summing over $i<j$,
$x$ and $x'$ with $y\ge g$ gives $\mathrm{Int}_g(k)$.

For~\eqref{eq:giant-split} note that
$\ones_{\{\mathcal G_g\ge1\}}=\mathcal G_g-(\mathcal G_g-1)^+$.
For~\eqref{eq:giant-crude} bound $T_\plus$ by $T_\plus(g)$ and $f_\plus$ by
$f^*_\plus(g)$ on the ranges of summation, and use~\eqref{eq:weight-masses}.
\end{proof}

\subsection{The boundary layer}\label{sec:boundary-proof}

We use two consequences of (RV) for the plus sign. First,
$T_\plus(n+1)/T_\plus(n)\to1$, so $f_\plus(n)/T_\plus(n)\to0$ and
\begin{equation}\label{eq:fstar-small}
 f^*_\plus(n)\le\sup_{m\ge n}\frac{f_\plus(m)}{T_\plus(m)}\;T_\plus(n)=o\bigl(T_\plus(n)\bigr).
\end{equation}
Second, a sum of $m'\ge1$ plus runs equal to $x$ has a summand of size at least
$x/m'$. Conditioning on the other summands gives
\begin{equation}\label{eq:one-summand}
 \Prob(S^\plus_{m'}=x)\le m'\,f^*_\plus(x/m').
\end{equation}

\begin{proof}[Proof of Theorem~\ref{thm:bigjump}(a)]
Fix $j\ge0$. If $j=0$ then $P_N(N)$ is $T_\plus(N)$ for a plus start and $0$ for a
minus start, which is $\kappa(0)T_\plus(N)$. Let $j\ge1$ and take
$g=g'_N=\lceil(N-j)/(j+1)\rceil$ in Lemma~\ref{lem:giant}. By
Lemma~\ref{lem:counting}, on $\{A_N=N-j\}$ some plus run has at least $g'_N$ slots,
so the first term of~\eqref{eq:giant-split} vanishes. Since $g'_N\ge N/(2j+2)$ for
large $N$, Lemma~\ref{lem:rv}(ii) gives $T_\plus(g'_N)\le C_jT_\plus(N)$.

The last term of~\eqref{eq:giant-split}. There are at most $j+1$ plus runs, so
$(\mathcal G_g-1)^+\le(j+1)\ones_{\{\mathcal G_g\ge2\}}$. The second part of
Lemma~\ref{lem:counting} gives
\[
 \E\bigl[(\mathcal G_g-1)^+\ones_{\{A_N=N-j\}}\bigr]\le(j+1)^3\,T_\plus(g'_N)^2
 =O\bigl(T_\plus(N)^2\bigr)=o\bigl(T_\plus(N)\bigr).
\]

The term $\mathrm{Int}_g$. It contains the configurations of the second sum
of~\eqref{eq:final-explicit}, a complete giant plus run followed by a final minus
run of length $m\le j$, which are the leading term of $P_N(N-j)$ when $\kappa(j)=0$.
By~\eqref{eq:giant-crude} and~\eqref{eq:fstar-small},
$\mathrm{Int}_g(N-j)\le\gamma^{(\sgn_0)}_\plus(j)f^*_\plus(g'_N)=o(T_\plus(g'_N))=o(T_\plus(N))$.

The term $\mathrm{Cen}_g$. Split the sum at $x=N/4$. For $x\le N/4$ and large $N$ we
have $N-j-x\ge N/2\ge g'_N$, and $T_\plus(N-j-x)/T_\plus(N)$ is at most
$T_\plus(\lfloor N/2\rfloor)/T_\plus(N)$, which is bounded by Lemma~\ref{lem:rv}(ii).
For each fixed $x$ the ratio tends to one. Since $\rho_j$ is a finite measure that
does not depend on $N$, dominated convergence and~\eqref{eq:weight-masses} give
\[
 \frac1{T_\plus(N)}\sum_{x\le N/4}\rho_j(x)\,T_\plus(N-j-x)\to\sum_{x\ge0}\rho_j(x)=\kappa(j).
\]
For $x>N/4$ we use~\eqref{eq:rho-rep}. The plus sums there have at most $j$
terms, and $\sum_mu^\minus_m(j)\le1$. So~\eqref{eq:one-summand} gives
$\rho_j(x)\le j\,f^*_\plus(N/(4j))$. Since $\sum_{x}T_\plus(N-j-x)\le\sum_{n\ge1}T_\plus(n)=\mu_\plus$,
\[
 \sum_{x>N/4}\rho_j(x)\,T_\plus(N-j-x)\le j\mu_\plus f^*_\plus\bigl(N/(4j)\bigr)=o\bigl(T_\plus(N)\bigr),
\]
by~\eqref{eq:fstar-small} and Lemma~\ref{lem:rv}(ii). Adding the four terms gives
$P_N(N-j)=\kappa(j)T_\plus(N)+o(T_\plus(N))$.
\end{proof}

At fixed $j$ the giant run is censored and its law enters only through $T_\plus$,
which is why (a) needs no local hypothesis.

\subsection{The big-jump zone}\label{sec:bigjump-proof}

In the zone we take $g=\lceil\eta L_k\rceil$ in Lemma~\ref{lem:giant} with
a small fixed $\eta$. Then~\eqref{eq:giant-split} has four terms, and we treat
them in three lemmas. Lemma~\ref{lem:no-giant} bounds the bins with no giant
run, Lemma~\ref{lem:two-giant} the bins with two, and Lemma~\ref{lem:one-giant}
evaluates $\mathrm{Cen}_g$ and $\mathrm{Int}_g$. Note that $k\ge L_k$ for
$pN<k\le N$, since $k-L_k=p(N-k)/(1-p)$.

We also need a bound for right deviations without big summands. It is the
elementary exponential estimate behind the inequalities of Fuk and
Nagaev~\cite{fuknagaev1971} and~\cite[Theorems~1.1 and~1.2]{nagaev1979}, and we
include the short proof.

\begin{lemma}[truncated right deviations]\label{lem:truncated}
Assume (RV) for the plus sign. For all $m\ge1$, $w\ge1$ and $z>0$,
\[
 \Prob\Bigl(S^\plus_m-m\mu_\plus\ge z,\ \max_{i\le m}X^\plus_i<w\Bigr)\le e^{C_{\mathrm K}}\bigl(m\,T_\plus(w)\bigr)^{z/(2w)} .
\]
\end{lemma}

\begin{proof}
If $mT_\plus(w)\ge1$ there is nothing to prove. Otherwise let
$\vartheta=\frac12\log(1/(mT_\plus(w)))>0$. For $0\le v\le\vartheta$ we have
$e^v\le1+v+\frac12v^2e^\vartheta$. We use this with $v=\vartheta X/w$ on
$\{X<w\}$, where $X=X_\plus$. By summation by parts and~\eqref{eq:karamata},
$\E[X^2\ones_{\{X<w\}}]\le2\sum_{n<w}nT_\plus(n)\le2C_{\mathrm K}w^2T_\plus(w)$. With
$1+v\le e^v$ this gives
$\E[e^{\vartheta X/w}\ones_{\{X<w\}}]\le\exp(\vartheta\mu_\plus/w+C_{\mathrm K}\vartheta^2e^\vartheta T_\plus(w))$.
Markov's inequality and independence give the bound
$\exp(-\vartheta z/w+C_{\mathrm K}m\vartheta^2e^\vartheta T_\plus(w))$. Here
$m\vartheta^2e^\vartheta T_\plus(w)=\vartheta^2e^{-\vartheta}\le1$, and
$e^{-\vartheta z/w}=(mT_\plus(w))^{z/(2w)}$.
\end{proof}

\begin{lemma}[no giant run]\label{lem:no-giant}
Assume (RV) and (SPAN). Let $0<\eta\le1/16$ and $r=1/(16\eta)$, so $r\ge1>1/\alpha_\plus$. There are
constants $C,c,C_8>0$ such that, for every fresh start, all large $N$ and all
$pN<k\le N$ with $L_k\ge C_8B_N$,
\[
 \Prob\bigl(A_N=k,\ M_N<\eta L_k\bigr)\le C\,\frac{(NT_\plus(L_k))^r}{b_{\plus,N}}
 +\frac{C}{b_{\minus,N}}\,e^{-cL_k/b_{\minus,N}}+C\,e^{-N^c}.
\]
\end{lemma}

\begin{proof}
Write $L=L_k$, $j=N-k$ and $P_0$ for the probability on the left. If $j=0$ then the
whole walk is one plus run and $P_0=0$. Let $j\ge1$, and take first a plus start.
On $\{A_N=k\}$ slot $N$ lies either in the $(m+1)$-th plus run, and then the first
$m$ minus runs sum to $j$ and $S^\plus_m<k\le S^\plus_{m+1}$, or in the $m$-th minus
run, and then $S^\plus_m=k$ and $S^\minus_{m-1}<j\le S^\minus_m$, as in the proof of
Proposition~\ref{prop:fresh-identity}. On $\{M_N<\eta L\}$ the plus runs that lie
inside the first $N$ slots are shorter than $\eta L$, and in the first case the final
run has $k-S^\plus_m<\eta L$ slots. The plus and minus runs are independent. Hence
$P_0\le E_1+E_2$ with
\begin{align*}
 E_1&=\sum_{m\ge0}u^\minus_m(j)\,q^\eta_{m+1},\qquad
 q^\eta_{m+1}=\sum_{1\le v<\eta L}\Prob\Bigl(S^\plus_m=k-v,\ \max_{i\le m}X^\plus_i<\eta L\Bigr)T_\plus(v),\\
 E_2&=\sum_{m\ge1}\Prob\Bigl(S^\plus_m=k,\ \max_{i\le m}X^\plus_i<\eta L\Bigr)\,q^\minus_m(j).
\end{align*}
For a minus start the plus run that contains slot $N$ is the $m$-th one, preceded by
$m$ minus runs and $m-1$ plus runs, and the minus run that contains slot $N$ is the
$(m+1)$-th one, preceded by $m$ plus runs. So $E_1$ becomes
$\sum_{m\ge1}u^\minus_m(j)q^\eta_m$, where $q^\eta_m$ is built on
$\Prob(S^\plus_{m-1}=k-v,\dots)$, and $E_2$ becomes
$\sum_{m\ge1}\Prob(S^\plus_m=k,\dots)\,q^\minus_{m+1}(j)$. We note below the two places
where this changes the argument.

We split the indices. Put $d_1=k-m\mu_\plus$ and $d_2=m\mu_\minus-j$. Then
$\mu_\minus d_1+\mu_\plus d_2=\mu_\minus k-\mu_\plus j=s(k-pN)=\mu_\minus L$. Let
$\mathcal A$ be the set of $m$ with $d_1\ge L/2$ and $\mathcal B$ its complement.
For $m\in\mathcal B$ we get $d_2\ge y_2=\mu_\minus L/(2\mu_\plus)$. The kernel sums
$\sum_mu^\minus_m(j)$, $\sum_mq^\minus_m(j)$, $\sum_mq^\minus_{m+1}(j)$,
$\sum_mq^\plus_m(k)$ and $\sum_mu^\plus_m(k)$ are at most $1$
by~\eqref{eq:kernel-sums}. So $E_1+E_2$ is at most the sum of the suprema over
$m\in\mathcal A$ of the plus factors and the suprema over $m\in\mathcal B$ of the
minus factors.

\emph{The indices in $\mathcal A$.} Let $m\in\mathcal A$ and let $v'$ be an
integer with $v'-m\mu_\plus\ge L/4$. Every $k-v$ in $q^\eta_{m+1}$ is such a $v'$,
since $k-v-m\mu_\plus\ge L/2-\eta L$, and so is $k$. For a minus start the
relevant sum $S^\plus_{m-1}$ has one term fewer, so its deviation $k-v-(m-1)\mu_\plus$
is larger by $\mu_\plus$, and the bound below applies with $m-1$ in place of $m$.
For $m\le1$ the probability $\Prob(S^\plus_m=v',\max_iX^\plus_i<\eta L)$ vanishes,
because $v'\ge L/4>\eta L$. For $m\ge2$ split the indices into two halves of sizes
$\lfloor m/2\rfloor$ and $\lceil m/2\rceil$. On $\{S^\plus_m=v'\}$ one half exceeds
its mean by at least $L/8$, and then the other half takes one exact value. By
Lemma~\ref{lem:truncated} with $w=\eta L$ and $z=L/8$, and by
Lemma~\ref{lem:stable}(ii),
\[
 \Prob\Bigl(S^\plus_m=v',\ \max_{i\le m}X^\plus_i<\eta L\Bigr)\le2e^{C_{\mathrm K}}\bigl(mT_\plus(\eta L)\bigr)^r\min\Bigl(1,\frac{\bar u_\plus}{a_\plus(\lfloor m/2\rfloor)}\Bigr).
\]
The indices in $\mathcal A$ satisfy $m\le m_+=k/\mu_\plus\le N/\mu_\plus$, and
$m_+\ge N/s$ since $k>pN$. Since $r>1/\alpha_\plus$, Lemma~\ref{lem:rv}(ii) for
$a_\plus$ shows that the right side is largest up to a constant at $m=m_+$. Also
$a_\plus(m_+)\ge a_\plus(N/s)=b_{\plus,N}$ and $T_\plus(\eta L)\le C_\eta T_\plus(L)$.
With $\sum_vT_\plus(v)\le1+\mu_\plus$ in $q^\eta_{m+1}$ the plus factors are at most
$C(NT_\plus(L))^r/b_{\plus,N}$.

\emph{The indices in $\mathcal B$.} Let $m\in\mathcal B$, so $x_m=m\mu_\minus-j\ge y_2$
and in particular $m\ge y_2/\mu_\minus$, which is large. On $\{S^\minus_m=j\}$ one half
of the minus runs falls below its mean by at least $x_m/2$, and the other half
takes one exact value. Lemma~\ref{lem:stable}(iv) for the half, with
$a_\minus(\lceil m/2\rceil)\le a_\minus(m)$, and Lemma~\ref{lem:stable}(ii) give
\[
 u^\minus_m(j)\le2e^{C_-}e^{-x_m/(2a_\minus(m))}\min\Bigl(1,\frac{\bar u_\minus}{a_\minus(\lfloor m/2\rfloor)}\Bigr).
\]
Next, $q^\minus_m(j)=\sum_{v\ge1}T_\minus(v)u^\minus_{m-1}(j-v)\le\mu_\minus\sup_{v\ge1}u^\minus_{m-1}(j-v)$, and
$(m-1)\mu_\minus-(j-v)\ge x_m-\mu_\minus\ge x_m/2$. So $q^\minus_m(j)$ obeys the same bound
with $\mu_\minus$ in front, $x_m/4$ in the exponent and $\lfloor(m-1)/2\rfloor$ in
place of $\lfloor m/2\rfloor$. For a minus start,
$q^\minus_{m+1}(j)=\sum_{v\ge1}T_\minus(v)u^\minus_m(j-v)\le\mu_\minus\sup_{v\ge1}u^\minus_m(j-v)$
and $m\mu_\minus-(j-v)\ge x_m$, so it obeys the bound for $u^\minus_m(j)$ with
$\mu_\minus$ in front. Let $m\le2N/\mu_\minus$. Then $t=a_\minus(m)\le c_5b_{\minus,N}$,
and Lemma~\ref{lem:rv}(ii) gives $a_\minus(m)\le c_6a_\minus(\lfloor(m-1)/2\rfloor)$. So
all these bounds are at most $Ct^{-1}e^{-y_2/(4t)}$. The function
$t\mapsto t^{-1}e^{-y_2/(4t)}$ increases on $(0,y_2/4]$. Choose
$C_8=8c_5\mu_\plus/\mu_\minus$. Then $L\ge C_8B_N$ gives $c_5b_{\minus,N}\le y_2/4$,
and the bounds are at most $Cb_{\minus,N}^{-1}e^{-cL/b_{\minus,N}}$. Let
$m>2N/\mu_\minus$. Then $x_m\ge m\mu_\minus-N\ge m\mu_\minus/2$, and
Lemma~\ref{lem:rv}(ii) gives $a_\minus(m)\le m^{1-2c}$ for large $m$ with
$c=(1-1/\alpha_\minus)/4$. So the bounds are at most $Ce^{-cm^{2c}}\le Ce^{-N^c}$.

Adding the bounds for $\mathcal A$ and $\mathcal B$ proves the lemma.
\end{proof}

The second term of Lemma~\ref{lem:no-giant}, from minus runs falling short of their
mean, is the reason for $\lambda_N$.

\begin{lemma}[two giant runs]\label{lem:two-giant}
Assume (RV) and (LR) for the plus sign, and let $0<\eta<1$. There is a constant
$C$ such that, for every fresh start, all $N$ and all $pN<k\le N$ with
$\eta L_k\ge1$ and $g=\lceil\eta L_k\rceil$,
\[
 \E\bigl[(\mathcal G_g-1)^+\ones_{\{A_N=k\}}\bigr]\le C\Bigl[(j+1)^2\,\frac{T_\plus(L_k)^2}{L_k}+(j+1)\,T_\plus(L_k)^2\Bigr].
\]
\end{lemma}

\begin{proof}
Since $(n-1)^+\le n(n-1)/2$ for integers $n\ge0$, it suffices to bound the
expected number of pairs of giant runs on $\{A_N=k\}$. The second run of a pair
is complete or final. As in the proof of Lemma~\ref{lem:giant}, the walk between
the two runs is a fresh minus walk until a plus run starts after $i_2\ge1$ of its
minus slots, with $x_2$ plus slots before it, which has probability
$\rho^{(\minus)}_{i_2}(x_2)$. For two complete runs of lengths $y_1,y_2$ the
expected number of pairs is
\[
 \sum_{i_1+i_2+i_3=j}\ \sum_{x_1,x_2,x_3}\rho^{(\sgn_0)}_{i_1}(x_1)\,\rho^{(\minus)}_{i_2}(x_2)\,\beta_{i_3}(x_3)
 \sum_{y_1+y_2=N-j-x_1-x_2-x_3,\ y_1,y_2\ge g}f_\plus(y_1)f_\plus(y_2),
\]
with $i_2,i_3\ge1$. For a complete run followed by the final run the expected
number is the same with $\beta_{i_3}(x_3)$ removed, $i_1+i_2=j$, and $T_\plus(y_2)$
in place of $f_\plus(y_2)$. By (LR), $f_\plus(n)\le C_{\mathrm{LR}}T_\plus(n)/n$ for all $n\ge1$,
so $f^*_\plus(g)\le C_{\mathrm{LR}}T_\plus(g)/g$. Let $y_1+y_2=y$ with both at least
$g$, so $y/2\ge g$. The larger of $y_1,y_2$ is at least $y/2$. So
\[
 \sum f_\plus(y_1)f_\plus(y_2)\le2f^*_\plus(g)T_\plus(g),\qquad
 \sum f_\plus(y_1)T_\plus(y_2)\le f^*_\plus(g)\sum_{n\ge g}T_\plus(n)+T_\plus(g)^2,
\]
and the second sum is at most $(C_{\mathrm K}C_{\mathrm{LR}}+1)T_\plus(g)^2$ by~\eqref{eq:karamata}.
The weights have total mass at most $(1+\mu_\plus)(j+1)^2$ in the first case and
$j+1$ in the second, by~\eqref{eq:weight-masses}. Finally
$T_\plus(g)\le C_\eta T_\plus(L_k)$ and $g\ge\eta L_k$.
\end{proof}

The terms $\mathrm{Cen}_g$ and $\mathrm{Int}_g$ need the spread of the weights. A
plus run that starts after the $i$-th minus slot is preceded by about
$\bar x_i=i\mu_\plus/\mu_\minus$ plus slots, and the fluctuation is of order $B_N$. For
$z>0$ put
\[
 \mathrm{sp}^\rho_i(z)=\sum_{x\,:\,|x-\bar x_i|>z}\rho^{(\sgn_0)}_i(x),\qquad
 \mathrm{sp}^\beta_i(z)=\sum_{x\,:\,|x-\bar x_i|>z}\beta_i(x).
\]

\begin{lemma}[spread of the weights]\label{lem:spread}
Assume (RV). For every $\varepsilon>0$ there are $Z$ and $N_0$ such that
$\mathrm{sp}^\rho_i(ZB_N)\le\varepsilon$ and $\mathrm{sp}^\beta_i(ZB_N)\le\varepsilon$ for
all $N\ge N_0$ and all $0\le i\le N$.
\end{lemma}

\begin{proof}
Let $\varpi_\sgn(z)=\sup_{m\ge1}\Prob(|S^\sgn_m-m\mu_\sgn|>z\,a_\sgn(m))$. By
Lemma~\ref{lem:stable}(i) the laws of $(S^\sgn_m-m\mu_\sgn)/a_\sgn(m)$ converge, so they are
tight, and $\varpi_\sgn(z)\to0$ as $z\to\infty$. By Lemma~\ref{lem:rv}(ii) there is $c_7$
with $a_\sgn(m)\le c_7B_N$ for $m\le4N/\mu_\minus$ and large $N$.

The number of minus runs. By (K3), $\{S^\minus_m=i\}\subset\{J_\minus(i)=m\}$, so for
$w>0$
\[
 \sum_{m\,:\,|m\mu_\minus-i|>w}u^\minus_m(i)\le\Prob\bigl(|J_\minus(i)\mu_\minus-i|>w\bigr)
 \le\Prob\bigl(S^\minus_{m_1}\ge i\bigr)+\Prob\bigl(S^\minus_{m_2}<i\bigr),
\]
with $m_1=\lfloor(i-w)/\mu_\minus\rfloor$ (the first term is absent if $i\le w$) and
$m_2=\lfloor(i+w)/\mu_\minus\rfloor$. Indeed $J_\minus(i)>(i+w)/\mu_\minus$ gives
$J_\minus(i)>m_2$, and then $S^\minus_{m_2}\le S^\minus_{J_\minus(i)-1}<i$. The first event is a deviation of at least $w$
above the mean, and the second a deviation of at least $w-\mu_\minus$ below it. For
$w\le N$ both indices are at most $3N/\mu_\minus$. So for $w=zB_N$ the right side is at
most $2\varpi_\minus(z/(2c_7))$ for large $N$.

The bound for $\mathrm{sp}^\rho_i$. Take a plus start. On $\{S^\minus_m=i\}$ we have $m\le i\le N$ and
$|S^\plus_m-\bar x_i|\le|S^\plus_m-m\mu_\plus|+(\mu_\plus/\mu_\minus)|m\mu_\minus-i|$.
By~\eqref{eq:rho-rep},
\begin{align*}
 \mathrm{sp}^\rho_i(ZB_N)&\le\sup_{m\le N}\Prob\bigl(|S^\plus_m-m\mu_\plus|>ZB_N/2\bigr)
 +\sum_{m\,:\,|m\mu_\minus-i|>\mu_\minus ZB_N/(2\mu_\plus)}u^\minus_m(i)\\
 &\le\varpi_\plus\Bigl(\frac Z{2c_7}\Bigr)+2\varpi_\minus\Bigl(\frac{\mu_\minus Z}{4\mu_\plus c_7}\Bigr).
\end{align*}
For a minus start~\eqref{eq:rho-rep} has $S^\plus_{m-1}$ in place of $S^\plus_m$, and
$|S^\plus_{m-1}-(m-1)\mu_\plus|>ZB_N/2-\mu_\plus$ is needed for the first event. For
large $N$ this is at least $ZB_N/4$, so the same bound holds with $Z/(4c_7)$ in place
of $Z/(2c_7)$ in the first term.

The bound for $\mathrm{sp}^\beta_i$. Write $S=S^\plus_{J-1}$ in~\eqref{eq:beta-rep}.
Then $\mathrm{sp}^\beta_i(z)$ is the expected number of integers $x\in[S,S+\xi_i]$ with
$|x-\bar x_i|>z$, which is at most
$(1+\xi_i)\ones_{\{|S-\bar x_i|>z/2\}}+\xi_i\ones_{\{\xi_i>z/2\}}$. Since $X^\plus_J$ is
independent of $(J,S)$ and $\xi_i\le X^\plus_J$,
\[
 \mathrm{sp}^\beta_i(z)\le(1+\mu_\plus)\Prob(|S-\bar x_i|>z/2)+\E\bigl[X_\plus\ones_{\{X_\plus>z/2\}}\bigr].
\]
The last term tends to zero as $z\to\infty$. For the first, $J$ is independent of the
plus runs, and
$\Prob(J>3N/\mu_\minus)\le\Prob(S^\minus_{\lfloor3N/\mu_\minus\rfloor}<N)\to0$ uniformly
in $i\le N$ by Lemma~\ref{lem:stable}(iv). On $\{J\le3N/\mu_\minus\}$ we bound
$|S-\bar x_i|$ as for $\mathrm{sp}^\rho_i$, with $J-1$ in place of $m$. The plus part
is controlled by $\varpi_\plus$ at the index $J-1$, and the number of minus runs by
$\varpi_\minus$ through the display above, since $J=J_\minus(i)$ and
$|(J-1)\mu_\minus-i|\le|J\mu_\minus-i|+\mu_\minus$. Taking $z=ZB_N$ with $Z$
large and then $N$ large proves the lemma.
\end{proof}

\begin{lemma}[one giant run]\label{lem:one-giant}
Assume (RV) and (SPAN), let the start be fresh with sign $\sgn_0$, let
$0<\eta\le1/2$, and put $g=\lceil\eta L_k\rceil$. Write $\kappa=\kappa^{(\sgn_0)}_\plus$
and $\gamma=\gamma^{(\sgn_0)}_\plus$. Then $\mathrm{Cen}_g(k)=0$ when $\kappa(N-k)=0$,
$\mathrm{Int}_g(k)=0$ when $\gamma(N-k)=0$, and
\[
 \lim_{K\to\infty}\ \limsup_{N\to\infty}\ \sup_{k\,:\,L_k\ge KB_N,\ \kappa(N-k)>0}
 \Bigl|\frac{\mathrm{Cen}_g(k)}{\kappa(N-k)\,T_\plus(\lfloor L_k\rfloor)}-1\Bigr|=0 .
\]
If in addition (LR) holds for the plus sign, then also
\[
 \lim_{K\to\infty}\ \limsup_{N\to\infty}\ \sup_{k\,:\,L_k\ge KB_N,\ \gamma(N-k)>0}
 \Bigl|\frac{\mathrm{Int}_g(k)}{\gamma(N-k)\,f_\plus(\lfloor L_k\rfloor)}-1\Bigr|=0 .
\]
\end{lemma}

\begin{proof}
The two vanishing statements follow from~\eqref{eq:giant-crude}.
By~\eqref{eq:renewal-thm} and~\eqref{eq:kappa}, $\kappa(i)\to1/\mu_\minus$, and
by~\eqref{eq:gamma-asymp}, $\gamma(i)/i\to s/\mu_\minus^2$. Hence
\[
 \kappa_*=\inf\{\kappa(i)\mid\kappa(i)>0\}>0,\qquad
 \gamma_*=\inf\{\gamma(i)/(i+1)\mid\gamma(i)>0\}>0 .
\]
Write $L=L_k$, $\hat L=\lfloor L\rfloor$ and $j=N-k$. Note that
$N-j-x=L-(x-\bar x_j)$ by~\eqref{eq:Lk}. Fix $\varepsilon>0$ and take
$Z,N_0$ such that Lemma~\ref{lem:spread} holds with $Z/2$ in place of $Z$. Since
$\mathrm{sp}^\rho_i(z)$ and $\mathrm{sp}^\beta_i(z)$ decrease in $z$, the same bounds
hold at $ZB_N$.

The term $\mathrm{Cen}_g$. Let $\kappa(j)>0$ and write
$\mathrm{Cen}_g(k)-\kappa(j)T_\plus(\hat L)=\sum_x\rho_j(x)[T_\plus(N-j-x)\ones_{\{N-j-x\ge g\}}-T_\plus(\hat L)]$.
If $|x-\bar x_j|\le ZB_N$, then $|N-j-x-\hat L|\le ZB_N+1\le(Z/K+1/L)L$.
Since $Z/K\le\eta/2$ for large $K$, the indicator is $1$, and by
Lemma~\ref{lem:rv}(i) the bracket is at most $\omega_T(Z/K)T_\plus(\hat L)$, where
$\omega_T(\delta)\to0$ as $\delta\to0$ once $N$ is large. For the other $x$ the
bracket is at most $T_\plus(g)+T_\plus(\hat L)\le C_\eta T_\plus(\hat L)$, and their
total weight is $\mathrm{sp}^\rho_j(ZB_N)\le\varepsilon\le(\varepsilon/\kappa_*)\kappa(j)$.
Hence
\[
 \bigl|\mathrm{Cen}_g(k)-\kappa(j)T_\plus(\hat L)\bigr|\le\bigl(\omega_T(Z/K)+C_\eta\varepsilon/\kappa_*\bigr)\kappa(j)T_\plus(\hat L).
\]

The term $\mathrm{Int}_g$. Let $\gamma(j)>0$ and define the measure $\chi_j$ by
\[
 \chi_j(z)=\sum_{i<j}\ \sum_{x+x'=z}\rho_i(x)\beta_{j-i}(x').
\]
It has total mass $\gamma(j)$ by~\eqref{eq:weight-masses}, and
\[
 \mathrm{Int}_g(k)=\sum_z\chi_j(z)f_\plus(N-j-z)\ones_{\{N-j-z\ge g\}} .
\]
If $|z-\bar x_j|>ZB_N$, then the split $(x,x')$ has $|x-\bar x_i|>ZB_N/2$ or
$|x'-\bar x_{j-i}|>ZB_N/2$. So Lemma~\ref{lem:spread} and~\eqref{eq:weight-masses}
bound the mass of these $z$ by
\[
 \sum_{i<j}\Bigl[(1+\mu_\plus)\mathrm{sp}^\rho_i(ZB_N/2)+\kappa(i)\mathrm{sp}^\beta_{j-i}(ZB_N/2)\Bigr]\le(2+\mu_\plus)\,j\,\varepsilon
 \le\frac{(2+\mu_\plus)\varepsilon}{\gamma_*}\,\gamma(j).
\]
Under (LR), $x\mapsto T_\plus(x)/x$ is regularly varying and
$f_\plus(n)\sim\alpha_\plus T_\plus(n)/n$. So Lemma~\ref{lem:rv}(i) gives
$|f_\plus(n)/f_\plus(\hat L)-1|\le\omega_f(Z/K)$ for $|n-\hat L|\le ZB_N+1$ and
large $N$, and Lemma~\ref{lem:rv}(ii) gives $f^*_\plus(g)\le C_\eta f_\plus(\hat L)$.
As for $\mathrm{Cen}_g$,
\[
 \bigl|\mathrm{Int}_g(k)-\gamma(j)f_\plus(\hat L)\bigr|\le\Bigl(\omega_f(Z/K)+\frac{(C_\eta+1)(2+\mu_\plus)\varepsilon}{\gamma_*}\Bigr)\gamma(j)f_\plus(\hat L).
\]
Let $N\to\infty$, then $K\to\infty$, then $\varepsilon\to0$.
\end{proof}

\begin{proof}[Proof of Theorem~\ref{thm:bigjump}(b) and (c)]
Write $\kappa,\gamma$ as in Lemma~\ref{lem:one-giant}, $L=L_k$,
$\hat L=\lfloor L\rfloor$ and $j=N-k$. Note that $T_\plus(\hat L)\ge T_\plus(L)$.
By~\eqref{eq:renewal-thm} and~\eqref{eq:gamma-asymp} there is $J_0\ge1$ with
$\kappa(j)\ge1/(2\mu_\minus)$ and $\gamma(j)\ge js/(2\mu_\minus^2)$ for $j\ge J_0$. Fix
\[
 \eta=\min\Bigl\{\frac1{32},\ \frac1{2J_0+2}\Bigr\},\qquad g=\lceil\eta L\rceil,
\]
so that $r=1/(16\eta)\ge2$. Fix $\varepsilon>0$. By~\eqref{eq:giant-split},
\begin{align*}
 \bigl|P_N(k)-G_N(k)\bigr|\le{}&\Prob(A_N=k,\ M_N<\eta L)+|\mathrm{Cen}_g(k)-\kappa(j)T_\plus(\hat L)|\\
 &+|\mathrm{Int}_g(k)-\gamma(j)f_\plus(\hat L)|+\E\bigl[(\mathcal G_g-1)^+\ones_{\{A_N=k\}}\bigr].
\end{align*}
In the zone $k\ge pN+K\lambda_NB_N$ we have $L\ge k-pN\ge K\lambda_NB_N\ge KB_N$. By
Lemma~\ref{lem:one-giant} the second and third terms are at most
$\varepsilon[\kappa(j)T_\plus(\hat L)+\gamma(j)f_\plus(\hat L)]=\varepsilon G_N(k)$ for
$K\ge K_1(\varepsilon)$ and $N\ge N_1(K)$. It remains to compare the first and last
terms with $G_N(k)$. We use~\eqref{eq:NT-bounds} for the plus sign.

\emph{Small $j$.} Let $j<J_0$. For large $N$ we have
$\eta L\le\eta N\le(N-j)/(j+1)$, so Lemma~\ref{lem:counting} gives
$\Prob(A_N=k,M_N<\eta L)=0$. If $G_N(k)=0$, then $\kappa(j)=\gamma(j)=0$, so
$\mathrm{Cen}_g=\mathrm{Int}_g=0$ and $\mathcal G_g=0$ almost surely on $\{A_N=k\}$, and
by~\eqref{eq:giant-split} $P_N(k)=0$. If $G_N(k)>0$, then
$G_N(k)\ge\min(\kappa_*T_\plus(\hat L),\gamma_*f_\plus(\hat L))\ge cT_\plus(L)/L$
by (LR). Also $L\ge N/2$ for large $N$. By Lemma~\ref{lem:two-giant} the last term is
at most $CJ_0^2T_\plus(L)^2$, and its ratio to $G_N(k)$ is at most
$CJ_0^2LT_\plus(L)$, which tends to zero since $\alpha_\plus>1$.

\emph{Large $j$.} Let $j\ge J_0$. By (LR), $f_\plus(\hat L)\ge\alpha_\plus T_\plus(\hat L)/(2L)$
for large $L$. Since $js/\mu_\minus^2=(N-L)/\mu_\minus$ by~\eqref{eq:Lk},
\[
 G_N(k)\ge\frac{T_\plus(L)}{2\mu_\minus}\Bigl(\frac{\alpha_\plus(N-L)}{2L}+1\Bigr)\ge c_8\,\frac{N}{L}\,T_\plus(L).
\]
By Lemma~\ref{lem:two-giant}, with $j+1\le N$, the last term divided by $G_N(k)$ is
at most $C(j+1)T_\plus(L)\le CNT_\plus(L)\le CK^{-(\alpha_\plus+1)/2}$
by~\eqref{eq:NT-bounds}, since $L/b_{\plus,N}\ge L/B_N\ge K$. By
Lemma~\ref{lem:no-giant}, for $K\ge C_8$, the first term divided by $G_N(k)$ is at
most $C$ times
\[
 \bigl(NT_\plus(L)\bigr)^{r-1}\frac{L}{b_{\plus,N}}
 +\frac{L}{b_{\minus,N}}\,e^{-cL/b_{\minus,N}}\,\frac1{NT_\plus(L)}
 +\frac{e^{-N^c}}{T_\plus(N)} .
\]
The last part tends to zero. The first part is at most
$C(L/b_{\plus,N})^{(1-\alpha_\plus)/2}$ by~\eqref{eq:NT-bounds} and $r\ge2$, and this
is at most $CK^{(1-\alpha_\plus)/2}$. For the middle
part put $x=L/b_{\minus,N}$ and $\varrho=b_{\minus,N}/b_{\plus,N}$, so
$L/b_{\plus,N}=x\varrho$. By~\eqref{eq:NT-bounds} the middle part is at most
$Cx^{\alpha_\plus+2}\varrho^{\alpha_\plus+1}e^{-cx}$. If $\varrho\le1$, then
$x\ge L/B_N\ge K$ and the bound is at most $C\sup_{x\ge K}x^{\alpha_\plus+2}e^{-cx}$. If
$\varrho>1$, then $B_N=b_{\minus,N}$, $\lambda_N=1+\log\varrho$ and $x\ge K(1+\log\varrho)$.
For $cK\ge2(\alpha_\plus+1)$ this gives
$\varrho^{\alpha_\plus+1}e^{-cx/2}\le e^{(\alpha_\plus+1)\log\varrho-cK(1+\log\varrho)/2}\le e^{-cK/2}$,
and the bound is at most $Ce^{-cK/2}\sup_xx^{\alpha_\plus+2}e^{-cx/2}$. In every case the
first term is at most $\varepsilon G_N(k)$ for $K$ large and then $N$ large. This proves
the ratio form of (b).

For the additive form note that $\gamma(j)\le(1+\mu_\plus)j\le(1+\mu_\plus)N$ by~\eqref{eq:gamma},
$(N-L)/\mu_\minus\le N/\mu_\minus$, $f_\plus(\hat L)\le2C_{\mathrm{LR}}T_\plus(\hat L)/L$ by (LR), and $\kappa\le1\le N/L$. Hence
\begin{equation}\label{eq:G-size}
 G_N(k)\le C_{\mathrm G}\widehat G_N(k),\qquad \overline G_N(k)\le C_{\mathrm G}\widehat G_N(k),
\end{equation}
with $C_{\mathrm G}$ depending only on the laws, and the same holds at every horizon $n$
and for both starting signs whenever $L^{(n)}_k\ge1$. The ratio form gives the additive form. Replacing
$\gamma(j)$ by $js/\mu_\minus^2=(N-L)/\mu_\minus$ changes $G_N(k)$ by
$|\gamma(j)-js/\mu_\minus^2|f_\plus(\hat L)\le2C_{\mathrm{LR}}|\gamma(j)-js/\mu_\minus^2|\,T_\plus(\hat L)/L$.
By~\eqref{eq:gamma-asymp} the factor $|\gamma(j)-js/\mu_\minus^2|$ is at most
$\varepsilon N$ for $j\ge J_\varepsilon$ and at most a constant $C_\varepsilon$ for
$j<J_\varepsilon$, and in the second case $L\ge N/2$. In both cases the change is at
most $(2C_{\mathrm{LR}}\varepsilon+4C_{\mathrm{LR}}C_\varepsilon/N)\widehat G_N(k)$. This proves (b).

For (c), let $\delta_N\to0$ and $0\le j\le\delta_NN$. Then $k-pN\ge(1-p)N/2$ for large
$N$, and $\lambda_NB_N=o(N)$ by Lemma~\ref{lem:rv}(v). So for every $K$ these bins lie
in the zone of (b) once $N$ is large, and (b) gives
$\sup_{j\le\delta_NN}|P_N(N-j)/G_N(N-j)-1|\to0$. Here
$L=N-js/\mu_\minus\ge N(1-\delta_Ns/\mu_\minus)$, so $T_\plus(\hat L)/T_\plus(N)\to1$
uniformly by Lemma~\ref{lem:rv}(i). By (LR) and $\gamma(j)\le(1+\mu_\plus)j$,
\[
 \frac{\gamma(j)f_\plus(\hat L)}{\kappa(j)T_\plus(N)}\le\frac{C(j+1)}{\kappa_*N}\le\frac{C(\delta_N+1/N)}{\kappa_*}\to0
\]
for $\kappa(j)>0$. Hence $G_N(N-j)/(\kappa(j)T_\plus(N))\to1$ uniformly over these $j$,
which proves the second statement of (c). For the first, the same bound gives
$\gamma(j)f_\plus(\hat L)\le C(\delta_N+1/N)T_\plus(N)$ for all $j\le\delta_NN$, so
$|G_N(N-j)-\kappa(j)T_\plus(N)|=o(T_\plus(N))$ uniformly. Also
$\widehat G_N(N-j)\le2T_\plus(\hat L)\le4T_\plus(N)$ there, and the additive form of (b) gives
$|P_N(N-j)-G_N(N-j)|=o(T_\plus(N))$ uniformly. This proves (c).
\end{proof}

\begin{corollary}[the lower bound]\label{cor:bigjump-lower}
Under the hypotheses of Theorem~\ref{thm:bigjump}(b),
\[
 \lim_{K\to\infty}\ \liminf_{N\to\infty}\ \inf_{pN+KB_N\le k\le N,\ G_N(k)>0}
 \frac{P_N(k)}{G_N(k)}\ge1 .
\]
That is, for every $\varepsilon>0$ there is $K$ such that
$P_N(k)\ge(1-\varepsilon)G_N(k)$ for all $pN+KB_N\le k\le N$ and all large $N$.
\end{corollary}

\begin{proof}
The first term of~\eqref{eq:giant-split} is nonnegative, and it is the only term
whose bound in the proof of (b) used $\lambda_N$. So
$P_N(k)\ge\mathrm{Cen}_g(k)+\mathrm{Int}_g(k)-\E[(\mathcal G_g-1)^+\ones_{\{A_N=k\}}]$,
with $J_0$, $\eta$ and $g$ as in that proof. For $pN+KB_N\le k\le N$ we have
$L_k\ge KB_N$. Lemma~\ref{lem:one-giant} needs only this, and it gives
$\mathrm{Cen}_g(k)+\mathrm{Int}_g(k)\ge(1-\varepsilon)G_N(k)$ for $K$ large and then $N$
large. The bounds for the last term in the proof of (b) also use only $L_k\ge KB_N$,
and they make it at most $\varepsilon G_N(k)$.
\end{proof}

\subsection{Consequences and an example}\label{sec:bigjump-consequences}

\begin{corollary}[the infinite density]\label{cor:infinite-density}
Assume (RV) and (SPAN) for both signs and (LR) for the plus sign. Fix $v\in(p,1)$
and put $\ell=(v-p)/(1-p)$. For a fresh start of either sign,
\[
 \frac{P_N(\lfloor vN\rfloor)}{T_\plus(N)}\longrightarrow
 I_\plus(v)=\frac1{\mu_\minus}\Bigl[\alpha_\plus(1-\ell)\ell^{-\alpha_\plus-1}+\ell^{-\alpha_\plus}\Bigr].
\]
The function $I_\plus$ is not integrable at $v\downarrow p$.
\end{corollary}

\begin{proof}
Let $k=\lfloor vN\rfloor$. Then $k-pN\sim(v-p)N$, which exceeds $K\lambda_NB_N$ for
every $K$ once $N$ is large, by Lemma~\ref{lem:rv}(v). So Theorem~\ref{thm:bigjump}(b)
gives $P_N(k)\sim G_N(k)$. Here $j=N-k\sim(1-v)N=(1-p)(1-\ell)N$ and $L_k\sim\ell N$.
By~\eqref{eq:gamma-asymp}, $\gamma^{(\sgn_0)}_\plus(j)\sim(1-\ell)N/\mu_\minus$. By (LR) and (RV),
$f_\plus(\lfloor L_k\rfloor)\sim\alpha_\plus\ell^{-\alpha_\plus-1}T_\plus(N)/N$ and
$T_\plus(\lfloor L_k\rfloor)\sim\ell^{-\alpha_\plus}T_\plus(N)$. By~\eqref{eq:renewal-thm},
$\kappa^{(\sgn_0)}_\plus(j)\to1/\mu_\minus$. Adding the two terms of~\eqref{eq:G} gives
$I_\plus(v)$. Near $v=p$ the first term is of order $(v-p)^{-\alpha_\plus-1}$, which
is not integrable.
\end{proof}

Rebenshtok, Denisov, H\"anggi and Barkai~\cite[Eqs.~(8) and~(10)]{rdhb2014} (see
also~\cite[Eqs.~(38) and~(52)]{rdhb2014b}) obtained for the two-state velocity model
an infinite density proportional to $|v|^{-1-\alpha}[\alpha-(\alpha-1)|v|]$, formally
and in continuous time. Since
$\alpha(1-\ell)\ell^{-\alpha-1}+\ell^{-\alpha}=\ell^{-\alpha-1}[\alpha-(\alpha-1)\ell]$,
the function $I_\plus$ has this shape. The complete giant run alone gives
$\ell^{-\alpha-1}[\alpha-\alpha\ell]$. The censored final run supplies the term
$\ell^{-\alpha}$, which changes the coefficient of the linear term from $-\alpha$ to
$-(\alpha-1)$. So their formula already contains the censored run. When $b_{\plus,N}=B_N$,
Theorem~\ref{thm:bulk}(c) on $|k-pN|\le KB_N$ and Theorem~\ref{thm:bigjump}(b) on
$k-pN\ge KB_N$ together give every bin on the plus side with relative error $o(1)$ in
the double limit. The matching between the stable bulk and the infinite density,
which is formal in the physics literature, is this statement.

\begin{remark}[the double limit]\label{rem:no-single-formula}
The bulk profile and the big-jump kernel are different expressions on their
common ground. Neither the bulk profile nor the big-jump kernel nor their sum has
relative error $o(1)$ uniformly over the plus side. Assume $b_{\plus,N}=B_N$ and
(LR) for the plus sign, and let $k=\lceil pN+KB_N\rceil$ with $K$ fixed. By Theorem~\ref{thm:bulk}(c),
$B_NP_N(k)/h_N(K)\to1$, and $h_N(K)$ is bounded below by Lemma~\ref{lem:profile}(d).
On the other hand, $j=N-k\sim(1-p)N$ and $L_k\sim KB_N/(1-p)$, so
by~\eqref{eq:gamma-asymp} and (LR)
\[
 B_NG_N(k)=\frac{\alpha_\plus}{s}\,\frac{NT_\plus(L_k)}{K}\,(1+o(1))=\alpha_\plus(1-p)^{\alpha_\plus}K^{-\alpha_\plus-1}(1+o(1)),
\]
where the second step uses $NT_\plus(b_{\plus,N})\to s$ and Lemma~\ref{lem:rv}(i), and
the censored term is negligible because $B_NT_\plus(L_k)\to0$. The right side is the
leading term of the tail of $h_N$, but it is not $h_N(K)$. Put
$r_N(K)=h_N(K)/(\alpha_\plus(1-p)^{\alpha_\plus}K^{-\alpha_\plus-1})$. Then
\[
 \frac{P_N(k)}{G_N(k)}=r_N(K)+o(1),\qquad
 \frac{P_N(k)}{B_N^{-1}h_N(K)+G_N(k)}=\frac{r_N(K)}{1+r_N(K)}+o(1).
\]
Since $h_N\le C_1$ by Lemma~\ref{lem:profile}(a), $r_N(K)$ is bounded in $N$ for each
$K$, so the second ratio stays below a constant smaller than one for large $N$. The
first ratio need not converge when $\theta_{\minus,N}$ oscillates, and it may equal
one at isolated values of $K$. In the example of Section~\ref{sec:verify}, where
$\theta_N=(1,1)$, $r_N(K)$ does not depend on $N$ and moves from $0.31$ at $K=1$ to
$1.68$ at $K=4$, so it equals one at some $K$ in between. So no fixed $K$ makes the
kernel accurate in general. Section~\ref{sec:verify} reports the check. By
Theorem~\ref{thm:bigjump}(b) and Theorem~\ref{thm:bulk}(c),
$\lim_{K\to\infty}\limsup_{N\to\infty}|r_N(K)-1|=0$, and the ratio $P_N/G_N$ tends to
one in this double limit, which is the form of Theorem~\ref{thm:bigjump}(b).
\end{remark}

\begin{remark}[the start of the zone]\label{rem:theta}
The zone of Theorem~\ref{thm:bigjump}(b) cannot start at $pN+KB_N$ with $K$ fixed
when $\theta_{\plus,N}\to0$. Some factor in front of $B_N$ that tends to infinity is
needed. Assume (LR) for the plus sign and let $\theta_{\plus,N}\to0$ along a
subsequence. Fix $K$ and let $k=\lceil pN+KB_N\rceil$. By Theorem~\ref{thm:bulk}(c)
and Lemma~\ref{lem:profile}(d), $B_NP_N(k)$ is bounded below. By (LR) and the bounds
$\gamma(j)\le(1+\mu_\plus)N$ and $\kappa\le1$,
\[
 B_NG_N(k)\le C\,NT_\plus(L_k)\,\frac{B_N}{L_k}+B_NT_\plus(L_k)
 \le\frac CK\,NT_\plus(L_k)+B_NT_\plus(B_N).
\]
Since $L_k\ge KB_N=Kb_{\plus,N}/\theta_{\plus,N}$, the bound~\eqref{eq:NT-bounds} gives
$NT_\plus(L_k)\le4s\,(K/\theta_{\plus,N})^{-(\alpha_\plus+1)/2}$, which tends to zero,
and $B_NT_\plus(B_N)\to0$. So $P_N(k)/G_N(k)\to\infty$ along the subsequence. The bins
there are carried by short minus runs, the second term of Lemma~\ref{lem:no-giant},
and not by a giant plus run. We do not show that the factor must be as large as
$\lambda_N$. A left-tail bound of order $\exp(-cy^\beta)$ in place of
Lemma~\ref{lem:stable}(iv) would replace $\lambda_N$ by $\lambda_N^{1/\beta}$, and we
do not pursue this. The lower bound needs no factor at all
(Corollary~\ref{cor:bigjump-lower}).
\end{remark}

\begin{example}[(RV) without (LR)]\label{ex:rv-not-lr}
Fix $1<\alpha<2$ and let $n_i=2^{2^i}$ for $i\ge1$ and $n_0=0$. Put
$h(n)=1/\log n_i$ for $n_{i-1}<n\le n_i$, and
\[
 T(n)=\frac{n^{-\alpha}\bigl(1+h(n)\bigr)}{1+h(1)},\qquad f(n)=T(n)-T(n+1)\qquad(n\ge1).
\]
Let both run laws be this law, so $p=1/2$ and $L_k=2k-N$. Then (RV) and (SPAN)
hold and (LR) fails. Let $N_i=2n_i$ and $k_i=3n_i/2$. Then $k_i-pN_i=n_i/2$, so $k_i$
lies in the zone of Theorem~\ref{thm:bigjump}(b) for every $K$ once $i$ is large,
and
\[
 \frac{P_{N_i}(k_i)}{G_{N_i}(k_i)}\to0 .
\]
So the conclusion of Theorem~\ref{thm:bigjump}(b), whose kernel reads $f_\plus$ at a
single point, fails under (RV) and (SPAN) alone, while (a) holds.
\end{example}

\begin{proof}[Verification]
Write $h_i=1/\log n_i=1/(2^i\log2)$, so $h_{i+1}=h_i/2$ and $h$ is nonincreasing
and tends to zero. Put $c=1/(1+h(1))$. Then $T(1)=1$, and $T$ is decreasing and tends to
zero, so $f$ is a probability mass function with $f(n)>0$ for all $n$. This gives
(SPAN), and $T(n)\sim cn^{-\alpha}$ gives (RV). For $n_{i-1}<n<n_i$ the points
$n,n+1$ lie in the same block, so $f(n)=c(1+h_i)(n^{-\alpha}-(n+1)^{-\alpha})\le2c\alpha n^{-\alpha-1}$.
At the block ends,
\[
 f(n_i)=c\Bigl[\bigl(n_i^{-\alpha}-(n_i+1)^{-\alpha}\bigr)(1+h_{i+1})+n_i^{-\alpha}\,\frac{h_i}2\Bigr],
\]
so $\frac c2h_in_i^{-\alpha}\le f(n_i)\le Ch_in_i^{-\alpha}$. Since $n_ih_i\to\infty$,
$n_if(n_i)/T(n_i)\to\infty$ and (LR) fails.

Let $N=N_i$, $k=k_i$, $j=n_i/2$, so $L_k=n_i$. By~\eqref{eq:gamma-asymp},
$\gamma(j)\ge n_i/(2\mu)$ for large $i$, where $\mu$ is the common mean. Hence
\[
 G_N(k)\ge\gamma(j)f(n_i)\ge\frac{c}{4\mu}\,h_in_i^{1-\alpha}.
\]
For the upper bound take $\eta=1/32$ and $g=\lceil n_i/32\rceil$ in
Lemma~\ref{lem:giant}, and drop the last term of~\eqref{eq:giant-split}.
Lemma~\ref{lem:no-giant} uses only (RV) and (SPAN). With $L_k=n_i$, $N=2n_i$ and
$r=2$ it bounds the first term by $Cn_i^{2-2\alpha}$ plus terms that decay faster
than every power of $n_i$, and $n_i^{2-2\alpha}=o(h_in_i^{1-\alpha})$.
By~\eqref{eq:giant-crude}, $\mathrm{Cen}_g(k)\le T(g)\le Cn_i^{-\alpha}$. For
$\mathrm{Int}_g$ note that $n_{i-1}=\sqrt{n_i}<n_i/32$ and $n_{i+1}=n_i^2>2n_i$ for large $i$.
So $n_i$ is the only block end in $[g,N]$, and $f(n)\le Cn_i^{-\alpha-1}$ for the
other $n$ there. With the measure $\chi_j$ of the proof of Lemma~\ref{lem:one-giant},
\[
 \mathrm{Int}_g(k)\le Cn_i^{-\alpha-1}\gamma(j)+f(n_i)\sup_z\chi_j(z)\le Cn_i^{-\alpha}+Ch_in_i^{-\alpha}\sup_z\chi_j(z).
\]
It remains to show $\sup_z\chi_j(z)=o(j)$. Since $\chi_j=\sum_{i'<j}\rho_{i'}*\beta_{j-i'}$
and the two factors have masses at most $1$ and $1+\mu$,
$\sup_z(\rho_{i'}*\beta_{j-i'})(z)\le\min\{(1+\mu)\sup_x\rho_{i'}(x),\ \sup_x\beta_{j-i'}(x)\}$.
One of $i'$ and $j-i'$ is at least $j/2$. So it suffices that
$\sup_x\rho_n(x)\to0$ and $\sup_x\beta_n(x)\to0$ as $n\to\infty$. Let $m_1=\lfloor n/(2\mu)\rfloor$.
By~\eqref{eq:rho-rep} and Lemma~\ref{lem:stable}(ii),
$\rho_n(x)\le\Prob(S_{m_1}\ge n)+\bar u/a(m_1-1)$, because a renewal at $n$ with at most
$m_1$ runs forces $S_{m_1}\ge n$. Both terms tend to zero, by the law of large numbers
and since $a(m)\to\infty$. By~\eqref{eq:beta-rep} and the independence of $X^\plus_J$,
$\beta_n(x)\le\sum_{r\ge0}T(r)\Prob(S_{J-1}=x-r)$ with $T(0)=1$, so
$\beta_n(x)\le(1+\mu)[\Prob(J\le m_1)+\bar u/a(m_1-1)]$, and
$\Prob(J\le m_1)\le\Prob(S_{m_1}\ge n)$, which tends to zero. Hence
$P_N(k)=o(h_in_i^{1-\alpha})$, and the ratio tends to zero.
\end{proof}

The spike at $n_i$ enters $G_N(k_i)$ at full weight but enters $P_N(k_i)$ only
through a window of width of order $B_N$ around $L_k$, as in the local big-jump
asymptotics of Denisov, Dieker and Shneer~\cite[Theorem~9.1]{dds2008}, which assume
a local hypothesis.

\begin{remark}[the hypothesis (LR)]\label{rem:wr}
The proof of Theorem~\ref{thm:bigjump}(b) uses (LR) to replace $f_\plus(n)$ by
$f_\plus(\lfloor L_k\rfloor)$ for $|n-L_k|\le ZB_N$, to bound $f^*_\plus(\eta L_k)$ by
$CT_\plus(L_k)/L_k$, and to bound $G_N$ below. The sharp condition is presumably a
regularity of $f_\plus$ over windows of width of order $B_N$. Example~\ref{ex:rv-not-lr}
shows that (RV) alone does not suffice. Whether some local condition strictly
weaker than (LR) suffices is open. At the junction
$L_k\asymp KB_N$ such windows have relative width of order $1/K$, and window
regularity at relative width $o(1)$ together with (RV) implies (LR), since
$T_\plus(n)-T_\plus((1+\delta)n)$ is then both $\delta nf_\plus(n)(1+o_\delta(1))$ and
$T_\plus(n)(\alpha_\plus\delta+O(\delta^2))$.
\end{remark}
